\chapter{Retriever Agent}

Il Retriever \`e l'agente pi\`u complesso: deve recuperare oggetti fisicamente sulla griglia,
gestire la capacit\`a di carico, navigare verso il magazzino, ricaricarsi e uscire, tutto
mantenendo una comunicazione attiva con Coordinator e altri Retriever.

\section{Gerarchia delle Priorit\`a}

Il ciclo decisionale del Retriever \`e una gerarchia di priorit\`a decrescenti, valutate
in sequenza ad ogni passo della fase \texttt{STAGE\_DECIDE}:

\begin{enumerate}
    \item \textbf{P1 — Consegna}: Se trasporta almeno un oggetto, si dirige al magazzino
          tramite la sequenza di consegna. Questa priorit\`a non \`e mai interrotta da
          nuovi compiti assegnati: il Retriever completa sempre la consegna prima di
          iniziare il prossimo recupero.
    \item \textbf{P2 — Ricarica d'emergenza}: Se l'energia scende sotto la soglia critica
          e non sta gi\`a consegnando, il Retriever interrompe qualsiasi operazione e si
          dirige al magazzino per ricaricarsi. La soglia \`e dinamica: se il magazzino
          pi\`u vicino noto dista $d$ celle, la soglia di allarme \`e $d \times 1.5$
          unit\`a di energia, garantendo sempre un margine di sicurezza.
    \item \textbf{P3 — Esecuzione compiti assegnati}: Se ha una coda di compiti assegnati
          dal Coordinator, preleva il prossimo oggetto dalla coda e si dirige verso di esso.
    \item \textbf{P4 — Auto-assegnazione (modalit\`a hive-mind)}: In assenza di compiti
          assegnati, il Retriever scansiona l'intero \texttt{known\_objects} accumulato
          (non solo le celle visibili) e si auto-assegna l'oggetto pi\`u conveniente
          tramite \texttt{try\_claim\_object} atomico. Questo meccanismo elimina i tempi
          morti tra un'assegnazione e l'altra, sfruttando la conoscenza condivisa tramite
          le mappe ricevute dagli altri agenti.
    \item \textbf{P5 — Mini-esplorazione}: Se non ha n\'e compiti n\'e oggetti noti,
          il Retriever esegue una breve frontier-based exploration locale per scoprire oggetti
          nell'area circostante.
\end{enumerate}

\section{Sequenza di Consegna al Magazzino}

La navigazione al magazzino avviene in quattro sotto-fasi gestite da una macchina a stati
interna (\texttt{\_wh\_step}):

\begin{enumerate}
    \item \textbf{Approach}: Il Retriever si dirige verso la cella \texttt{WAREHOUSE\_ENTRANCE}
          pi\`u vicina del magazzino ottimale (selezionato da \texttt{get\_best\_warehouse\_for}).
          Il percorso A* usa le forbidden types per forzare il passaggio dall'ingresso,
          escludendo le celle di uscita come nodi intermedi.
    \item \textbf{Deposit}: Entrato nel magazzino attraverso l'ingresso, il Retriever
          deposita tutti gli oggetti trasportati, che vengono conteggiati nelle metriche
          e rimossi dalla griglia. Lo stato passa a \texttt{RECHARGE} se l'energia \`e
          sotto soglia, altrimenti direttamente a \texttt{EXIT}.
    \item \textbf{Recharge}: Il Retriever rimane sulla cella magazzino mentre viene ricaricato
          a tasso costante (\texttt{recharge\_rate} unit\`a per passo) fino al recupero
          completo dell'energia massima. Durante la ricarica continua a ricevere messaggi.
    \item \textbf{Exit}: Il Retriever si dirige verso la cella \texttt{WAREHOUSE\_EXIT} e
          la attraversa. Uscito dal magazzino, \texttt{\_wh\_step} viene azzerato; prima
          di tornare all'esplorazione, il Retriever tenta immediatamente una nuova
          auto-assegnazione tramite \texttt{known\_objects}, riducendo i passi a vuoto
          tra una consegna e la successiva.
\end{enumerate}

\section{Selezione del Magazzino Ottimale}

La funzione \texttt{get\_best\_warehouse\_for} seleziona la stazione di deposito/ricarica
ottimale tra tutte quelle note all'agente. Il punteggio per ogni stazione combina tre fattori:

\begin{itemize}
    \item \textbf{Distanza}: La distanza euclidea dalla posizione corrente del Retriever.
    \item \textbf{Capacit\`a}: Le stazioni con pi\`u Retriever gi\`a in coda vengono
          penalizzate per distribuire il carico.
    \item \textbf{Fattibilit\`a energetica}: Se la distanza stimata (moltiplicata per un
          fattore di percorso 1.4 per tenere conto degli ostacoli) supera l'energia residua
          dell'agente, la stazione riceve una penalizzazione aggiuntiva di 200 punti,
          rendendola di fatto non preferibile. Questo impedisce al Retriever di dirigersi
          verso un magazzino irraggiungibile e rimanere senza energia a met\`a percorso.
\end{itemize}

\section{Peer Negotiation tra Retriever}

I conflitti tra Retriever che puntano allo stesso oggetto si risolvono in due livelli:

\begin{enumerate}
    \item \textbf{Claim atomico}: Ogni auto-assegnazione passa per
          \texttt{CommunicationManager.try\_claim\_object}, operazione atomica condivisa
          da tutti gli agenti. Solo un Retriever ottiene il claim; i contendenti ricevono
          immediatamente un fallimento e selezionano un altro oggetto.
    \item \textbf{Aggiornamento reattivo via TaskStatusMessage}: Dopo ogni claim
          riuscito, il Retriever diffonde un \texttt{TaskStatusMessage} a tutti i vicini.
          Alla ricezione, ogni peer rimuove dalla propria \texttt{task\_queue} gli oggetti
          gi\`a presenti nella coda del mittente (\texttt{PEER-YIELD}), aggiornando
          preventivamente la propria pianificazione senza attendere un fallimento esplicito.
\end{enumerate}

L'assenza di un arbitro centralizzato e la natura broadcast del \texttt{TaskStatusMessage}
garantiscono convergenza rapida: nella pratica, le collisioni residue dopo il claim atomico
sono rare e vengono risolte entro uno-due passi di comunicazione.
